\chapter{Algebraic and spectral-sequence background}
\label{chap:background}

The paper uses elementwise pages, cycles, boundaries, detection, and several
changes of grading.  Mathlib supplies categorical complexes and a generic
spectral-sequence record, but not this Adams-specific elementwise layer.  This
chapter separates the exact Mathlib roots from the project interfaces that
must still be constructed.

\section{Exact Mathlib roots}

% SOURCE: Mathlib/Data/ZMod/Defs.lean:142, declaration ZMod; Mathlib/Algebra/Field/ZMod.lean:29-30, prime-field instance.
\begin{definition}[The field $\mathbb F_2$]
  \label{def:mathlib-f2}
  \lean{ZMod}
  \mathlibok
  The coefficient field used throughout is
  $\mathbb F_2:=\operatorname{ZMod}(2)$ with its pinned Mathlib field
  structure.  The interface records $1+1=0$, characteristic two, decidable
  equality, and finiteness; later sign simplifications must cite these
  properties rather than treating every abelian category as characteristic
  two.
\end{definition}

% SOURCE: Mathlib/Algebra/Category/ModuleCat/Basic.lean, declaration ModuleCat.
\begin{definition}[$\mathbb F_2$-module category]
  \label{def:mathlib-f2-module-category}
  \lean{ModuleCat}
  \mathlibok
  Mathlib's category of modules over $\mathbb F_2=\operatorname{ZMod}(2)$ is
  the concrete abelian category in which the elementwise spectral-sequence
  adapter is constructed.
\end{definition}

% SOURCE: Mathlib/Algebra/Category/ModuleCat/Abelian.lean:75, instance ModuleCat.abelian.
\begin{theorem}[The module category is abelian]
  \label{thm:mathlib-module-category-abelian}
  \lean{ModuleCat.abelian}
  \mathlibok
  \uses{def:mathlib-f2,def:mathlib-f2-module-category}
  The category of modules over $\mathbb F_2$ carries Mathlib's abelian-category
  instance.
\end{theorem}

% SOURCE: Mathlib/CategoryTheory/Abelian/Basic.lean, declaration CategoryTheory.Abelian.
\begin{definition}[Abelian category]
  \label{def:mathlib-abelian}
  \lean{CategoryTheory.Abelian}
  \mathlibok
  In the pinned Mathlib declaration, an abelian category is preadditive, has
  kernels, cokernels, and finite products, and every monomorphism and
  epimorphism is normal.  A zero object and the usual finite biproduct
  consequences are derived from these fields.
\end{definition}

% SOURCE: Mathlib/Algebra/Homology/HomologicalComplex.lean, declaration CategoryTheory.HomologicalComplex.
\begin{definition}[Homological complex]
  \label{def:mathlib-homological-complex}
  \lean{HomologicalComplex}
  \mathlibok
  For a category $\mathcal C$, an index type $\iota$, and a complex shape
  $c$, a homological complex consists of objects $K_i$, differentials
  $d_{ij}:K_i\to K_j$ allowed by $c$, and the square-zero relations required
  by that shape.
\end{definition}

% SOURCE: Mathlib/Algebra/Homology/HomologicalComplex.lean, abbreviation CategoryTheory.ChainComplex.
\begin{definition}[Chain complex]
  \label{def:mathlib-chain-complex}
  \lean{ChainComplex}
  \mathlibok
  A chain complex in $\mathcal C$ indexed by an additive index type is the
  Mathlib homological complex with differential of degree $-1$.
\end{definition}

% SOURCE: Mathlib/Algebra/Homology/ShortComplex/Exact.lean, declaration CategoryTheory.ShortComplex.Exact.
\begin{definition}[Exact short complex]
  \label{def:mathlib-short-complex-exact}
  \lean{CategoryTheory.ShortComplex.Exact}
  \mathlibok
  For a short complex $X_1\xrightarrow{f}X_2\xrightarrow{g}X_3$, Mathlib
  exactness means that there exists homology data whose homology object is
  zero.  In an abelian category, the next node identifies this with the
  familiar image-to-kernel comparison being an isomorphism.
\end{definition}

% SOURCE: Mathlib/CategoryTheory/Abelian/Exact.lean:77, theorem CategoryTheory.ShortComplex.exact_iff_isIso_imageToKernel.
\begin{theorem}[Exactness via image and kernel]
  \label{thm:mathlib-exact-iff-image-kernel}
  \lean{CategoryTheory.ShortComplex.exact_iff_isIso_imageToKernel}
  \mathlibok
  \uses{def:mathlib-short-complex-exact,def:mathlib-abelian}
  In an abelian category, a short complex is exact if and only if its canonical
  morphism from the image of the first map to the kernel of the second map is
  an isomorphism.
\end{theorem}

% SOURCE: Mathlib/CategoryTheory/GradedObject.lean, declaration CategoryTheory.GradedObject.comap.
\begin{definition}[Regrading a Mathlib graded object]
  \label{def:mathlib-graded-comap}
  \lean{CategoryTheory.GradedObject.comap}
  \mathlibok
  A function $h:J\to I$ induces Mathlib's regrading functor from
  $I$-graded objects to $J$-graded objects by precomposition.  This declaration
  regrades objects only; it does not by itself regrade complexes or spectral
  sequences.
\end{definition}

% SOURCE: Mathlib/Algebra/Homology/ComplexShape.lean, declaration ComplexShape.up'.
\begin{definition}[Complex shape of a fixed additive degree]
  \label{def:mathlib-complex-shape-up}
  \lean{ComplexShape.up'}
  \mathlibok
  For an additive index type and a fixed degree $a$, Mathlib's
  $\operatorname{up}'(a)$ permits a differential from $i$ to $j$ exactly when
  $i+a=j$.  The AIM Adams shape therefore uses the integer pair
  $(r,r-1)$, not Mathlib's prepackaged cohomological convention $(r,1-r)$.
\end{definition}

% SOURCE: Mathlib/Algebra/Homology/SpectralSequence/Basic.lean, declaration CategoryTheory.SpectralSequence.Hom.
\begin{definition}[Mathlib spectral-sequence morphism]
  \label{def:mathlib-spectral-sequence-hom}
  \lean{CategoryTheory.SpectralSequence.Hom}
  \mathlibok
  A Mathlib morphism between spectral sequences has pagewise chain maps
  compatible with passage to homology, but both sequences must have the same
  index type, complex shape, and starting page.  It therefore cannot express
  the classical--synthetic reindexing required later.
\end{definition}

% SOURCE: Mathlib/Algebra/Homology/SpectralObject/Basic.lean, declaration CategoryTheory.Abelian.SpectralObject.
\begin{definition}[Mathlib abelian spectral object]
  \label{def:mathlib-spectral-object}
  \lean{CategoryTheory.Abelian.SpectralObject}
  \mathlibok
  A Mathlib spectral object in an abelian category is an indexed family of
  homology functors equipped with functorial connecting maps and the three
  exactness conditions forming its long exact sequences.
\end{definition}

% SOURCE: Mathlib/Algebra/Homology/SpectralObject/HasSpectralSequence.lean:67, declaration CategoryTheory.Abelian.SpectralObject.SpectralSequenceDataCore.
\begin{definition}[Mathlib spectral-sequence recipe]
  \label{def:mathlib-spectral-sequence-data-core}
  \lean{CategoryTheory.Abelian.SpectralObject.SpectralSequenceDataCore}
  \mathlibok
  \uses{def:mathlib-spectral-object,def:mathlib-complex-shape-up}
  A spectral-sequence data core records the grading degree and four monotone
  spectral-object indices that produce each page object and its differential
  shape.
\end{definition}

% SOURCE: Mathlib/Algebra/Homology/SpectralObject/HasSpectralSequence.lean:303, declaration CategoryTheory.Abelian.SpectralObject.HasSpectralSequence.
\begin{definition}[Mathlib spectral-object admissibility]
  \label{def:mathlib-has-spectral-sequence}
  \lean{CategoryTheory.Abelian.SpectralObject.HasSpectralSequence}
  \mathlibok
  \uses{def:mathlib-spectral-object,def:mathlib-spectral-sequence-data-core}
  For a spectral object and a page recipe, Mathlib's
  \texttt{HasSpectralSequence} class contains the endpoint vanishing needed to
  identify page homology with the successor page.
\end{definition}

% SOURCE: Mathlib/Algebra/Homology/SpectralObject/SpectralSequence.lean:481, declaration CategoryTheory.Abelian.SpectralObject.spectralSequence.
\begin{definition}[Mathlib spectral sequence of a spectral object]
  \label{def:mathlib-spectral-object-spectral-sequence}
  \lean{CategoryTheory.Abelian.SpectralObject.spectralSequence}
  \mathlibok
  \uses{def:mathlib-spectral-object,def:mathlib-spectral-sequence-data-core,def:mathlib-has-spectral-sequence,def:mathlib-spectral-sequence}
  Given a spectral object, a data core, and the endpoint-vanishing instance,
  Mathlib constructs a \texttt{CategoryTheory.SpectralSequence} with the
  prescribed page shapes.
\end{definition}

% SOURCE: Mathlib/Algebra/Homology/SpectralObject/SpectralSequence.lean:650, declaration CategoryTheory.Abelian.SpectralObject.E₂SpectralSequence.
\begin{definition}[Mathlib $E_2$ spectral sequence]
  \label{def:mathlib-e2-spectral-sequence}
  \lean{CategoryTheory.Abelian.SpectralObject.E₂SpectralSequence}
  \mathlibok
  \uses{def:mathlib-spectral-object-spectral-sequence}
  A spectral object indexed by extended integers has Mathlib's canonical
  cohomological $E_2$ spectral sequence on $\mathbb Z\times\mathbb Z$.
\end{definition}

% SOURCE: Mathlib/Algebra/Homology/SpectralObject/Cycles.lean:58, declaration CategoryTheory.Abelian.SpectralObject.cycles.
\begin{definition}[Mathlib spectral-object cycles]
  \label{def:mathlib-spectral-object-cycles}
  \lean{CategoryTheory.Abelian.SpectralObject.cycles}
  \mathlibok
  \uses{def:mathlib-spectral-object}
  Mathlib's cycle object for a spectral-object interval is the appropriate
  kernel of its connecting morphism.
\end{definition}

% SOURCE: Mathlib/Algebra/Homology/SpectralObject/Cycles.lean:62, declaration CategoryTheory.Abelian.SpectralObject.opcycles.
\begin{definition}[Mathlib spectral-object opcycles]
  \label{def:mathlib-spectral-object-opcycles}
  \lean{CategoryTheory.Abelian.SpectralObject.opcycles}
  \mathlibok
  \uses{def:mathlib-spectral-object}
  Mathlib's opposite-cycle object is the corresponding cokernel.  It is
  categorical infrastructure, not yet the paper's chosen element subgroup
  $B_r$.
\end{definition}

% SOURCE: Mathlib/Algebra/Homology/HomotopyCategory/SpectralObject.lean:47, declaration HomotopyCategory.spectralObjectMappingCone.
\begin{definition}[Mapping-cone spectral object]
  \label{def:mathlib-mapping-cone-spectral-object}
  \lean{HomotopyCategory.spectralObjectMappingCone}
  \mathlibok
  \uses{def:mathlib-spectral-object,def:mathlib-chain-complex}
  Mathlib's mapping-cone construction is a spectral object valued in the
  homotopy category of cochain complexes.  Precomposing it with a filtered
  complex regarded as a filtration-indexed functor is the intended project
  route to a filtered-complex spectral sequence.
\end{definition}

% SOURCE: Mathlib/Algebra/Homology/DerivedCategory/Ext/Basic.lean, declaration CategoryTheory.Abelian.Ext.
\begin{definition}[Mathlib derived Ext]
  \label{def:mathlib-derived-ext}
  \lean{CategoryTheory.Abelian.Ext}
  \mathlibok
  In an abelian category satisfying Mathlib's Ext-size hypothesis,
  $\operatorname{Ext}(X,Y,n)$ is a singly graded derived Ext type indexed by
  $n\in\mathbb N$.  It does not provide the Steenrod algebra, internal grading,
  comodules, or the bigraded Adams $\operatorname{Ext}_{\mathcal A}^{s,t}$.
\end{definition}

% SOURCE: Mathlib/CategoryTheory/Shift/Basic.lean, declaration CategoryTheory.HasShift.
\begin{definition}[Mathlib shift structure]
  \label{def:mathlib-shift}
  \lean{CategoryTheory.HasShift}
  \mathlibok
  Mathlib's shift structure supplies a coherent family of autoequivalences
  indexed here by $\mathbb Z$.
\end{definition}

% SOURCE: Mathlib/CategoryTheory/Triangulated/Basic.lean, declaration CategoryTheory.Pretriangulated.Triangle.
\begin{definition}[Mathlib triangle]
  \label{def:mathlib-triangle}
  \lean{CategoryTheory.Pretriangulated.Triangle}
  \mathlibok
  A Mathlib triangle consists of $X\to Y\to Z\to X[1]$ in a category with
  integer shifts.
\end{definition}

% SOURCE: Mathlib/CategoryTheory/Triangulated/Pretriangulated.lean, declaration CategoryTheory.Pretriangulated.
\begin{definition}[Mathlib pretriangulated structure]
  \label{def:mathlib-pretriangulated}
  \lean{CategoryTheory.Pretriangulated}
  \mathlibok
  Mathlib's pretriangulated class chooses distinguished triangles and their
  closure laws.  It does not include a stable-spectrum model or compatibility
  of a monoidal product with distinguished triangles.
\end{definition}

% SOURCE: Mathlib/CategoryTheory/Triangulated/HomologicalFunctor.lean:82, declaration CategoryTheory.Functor.IsHomological.
\begin{definition}[Mathlib homological functor]
  \label{def:mathlib-homological-functor}
  \lean{CategoryTheory.Functor.IsHomological}
  \mathlibok
  \uses{def:mathlib-pretriangulated,def:mathlib-abelian}
  A functor from a pretriangulated category to an abelian category is
  homological when it preserves zero morphisms and sends every distinguished
  triangle to an exact short complex.
\end{definition}

% SOURCE: Mathlib/CategoryTheory/Triangulated/HomologicalFunctor.lean:217, declaration CategoryTheory.Functor.homologySequence_exact₂.
\begin{theorem}[Middle exactness of a homology sequence]
  \label{thm:mathlib-homology-sequence-exact-two}
  \lean{CategoryTheory.Functor.homologySequence_exact₂}
  \mathlibok
  \uses{def:mathlib-homological-functor,def:mathlib-triangle}
  A Mathlib homological functor sends the middle two maps of every shifted
  distinguished triangle to an exact short complex.
\end{theorem}

% SOURCE: Mathlib/CategoryTheory/Triangulated/HomologicalFunctor.lean:226, declaration CategoryTheory.Functor.homologySequence_exact₃.
\begin{theorem}[Connecting-map exactness of a homology sequence]
  \label{thm:mathlib-homology-sequence-exact-three}
  \lean{CategoryTheory.Functor.homologySequence_exact₃}
  \mathlibok
  \uses{def:mathlib-homological-functor,def:mathlib-triangle}
  Rotation gives exactness at the term adjacent to the homology-sequence
  connecting morphism.
\end{theorem}

% SOURCE: Mathlib/CategoryTheory/Triangulated/HomologicalFunctor.lean:234, declaration CategoryTheory.Functor.homologySequence_exact₁.
\begin{theorem}[Shifted exactness of a homology sequence]
  \label{thm:mathlib-homology-sequence-exact-one}
  \lean{CategoryTheory.Functor.homologySequence_exact₁}
  \mathlibok
  \uses{def:mathlib-homological-functor,def:mathlib-triangle}
  Inverse rotation gives the remaining adjacent exactness statement, allowing
  the three short complexes to splice into a long exact sequence.
\end{theorem}

% SOURCE: Mathlib/CategoryTheory/Monoidal/Category.lean, declaration CategoryTheory.MonoidalCategory.
\begin{definition}[Mathlib monoidal category]
  \label{def:mathlib-monoidal-category}
  \lean{CategoryTheory.MonoidalCategory}
  \mathlibok
  Mathlib's monoidal category provides tensor product, unit, associator, and
  unitors with their coherence laws.
\end{definition}

% SOURCE: Mathlib/CategoryTheory/Monoidal/Braided/Basic.lean, declaration CategoryTheory.SymmetricCategory.
\begin{definition}[Mathlib symmetric category]
  \label{def:mathlib-symmetric-category}
  \lean{CategoryTheory.SymmetricCategory}
  \mathlibok
  A symmetric category extends a monoidal category with a symmetric braiding.
  Exactness of tensoring a triangle remains extra project data.
\end{definition}

\section{Filtration completeness and page shapes}

% SOURCE: aimpaper/main.tex:253-271 and 827-862; strong convergence is assumed throughout.
\begin{definition}[Complete and separated decreasing filtration]
  \label{def:complete-separated-filtration}
  \uses{def:core-filtration,def:core-exhaustive-filtration}
  \notready
  For a decreasing filtration $F$ on a graded abelian group $A$, exhaustive
  means $\bigcup_{s\in\mathbb Z}F^sA=A$, separated means
  $\bigcap_{s\in\mathbb Z}F^sA=0$, and complete means that the canonical map
  \[
    A\longrightarrow\varprojlim_s A/F^sA
  \]
  is an isomorphism, degree by degree.  A convergence witness also records the
  regularity or boundedness condition that makes only finitely many filtration
  pieces relevant in each total degree.
\end{definition}

% SOURCE: aimpaper/main.tex:134-150, 264-295, 799-823, and 1277-1317.
\begin{definition}[Spectral-sequence page-level convention]
  \label{def:ss-page-level}
  \notready
  A page-level convention records the first page, successor operation, page
  index used in $Z_r/B_r$, and the translation from any auxiliary quotient
  level to that page.  In particular, a classical $E_r$ page corresponds to
  the synthetic quotient $S/\lambda^{r-1}$; the two integers must never be
  identified definitionally.
\end{definition}

% SOURCE: aimpaper/main.tex:134-150 and 799-807, Adams differential degree.
\begin{definition}[Classical and synthetic Adams complex shape]
  \label{def:adams-complex-shape}
  \uses{def:mathlib-complex-shape-up,def:ss-page-level}
  \notready
  On the bidegree index $\mathbb Z\times\mathbb Z$, the page-$r$ Adams
  differential has additive degree $(r,r-1)$, hence is represented by
  Mathlib's fixed-degree shape $\operatorname{up}'(r,r-1)$.  The synthetic
  shape uses the same first two coordinates and fixes the weight coordinate.
\end{definition}

% SOURCE: aimpaper/main.tex:264-295 and 999-1007, extension differential degree.
\begin{definition}[Extension spectral-sequence complex shape]
  \label{def:ess-complex-shape}
  \uses{def:mathlib-complex-shape-up,def:ss-page-level}
  \notready
  On the paper's $(s,t)$ indexing, an extension differential $d_n^f$ has
  additive degree $(n,n)$.  The corresponding Mathlib complex shape is
  specified pagewise by $\operatorname{up}'(n,n)$, with any spectrum side or
  synthetic weight retained as an additional fixed index.
\end{definition}

\section{Gradings and page presentations}

% SOURCE: aimpaper/main.tex:134-150 and 253-295.
\begin{definition}[Adams bidegree convention]
  \label{def:adams-bidegree}
  \notready
  A classical Adams class has bidegree $(s,t)$, Adams filtration $s$, and stem
  $t-s$.  A classical differential has
  \[
    d_r:E_r^{s,t}\longrightarrow E_r^{s+r,t+r-1},
  \]
  hence lowers the stem by one.  In an extension spectral sequence the page
  differential is instead written with the paper's filtration shift
  $(n,n)$; every comparison must record which convention it uses.
\end{definition}

% SOURCE: aimpaper/main.tex:690-739, 799-823, and 1138-1150.
\begin{definition}[Synthetic tridegree and page-level conversion]
  \label{def:synthetic-tridegree}
  \uses{def:adams-bidegree}
  \notready
  A synthetic Adams class is indexed by $(s,t,w)$, where $s$ is Adams
  filtration, $t-s$ is the topological stem, and $w$ is synthetic weight.  A
  synthetic $d_r$ has degree $(r,r-1,0)$.  Multiplication by $\lambda$ has
  degree $(0,0,-1)$.  The classical page $E_r$ corresponds to the quotient
  level $S/\lambda^{r-1}$; this translation is data, not a definitional
  equality.
\end{definition}

% SOURCE: aimpaper/main.tex:827-862, definitions of Z_r and B_r.
\begin{definition}[Chosen cycles, boundaries, and pages]
  \label{def:chosen-page-presentation}
  \uses{def:mathlib-spectral-sequence,def:adams-bidegree,def:ss-page-level,def:adams-complex-shape}
  \notready
  An elementwise presentation of a spectral sequence chooses nested subgroups
  $B_r^{s,t}\subseteq Z_r^{s,t}$ and identifications
  $E_r^{s,t}\cong Z_r^{s,t}/B_r^{s,t}$, with
  $Z_{r+1}$ the classes on which $d_r$ vanishes and $B_{r+1}$ enlarged by the
  image of $d_r$.  It also includes
  $Z_\infty=\bigcap_r Z_r$ and $B_\infty=\bigcup_r B_r$ when these operations
  exist in the chosen algebraic setting.
\end{definition}

% SOURCE: aimpaper/main.tex:134-151, 827-862, and throughout the appendix.
\begin{definition}[Cycle, boundary, and permanent cycle]
  \label{def:ss-cycle-boundary-permanent}
  \uses{def:chosen-page-presentation}
  \notready
  A representative is an $r$-cycle when it lies in $Z_r$, an $r$-boundary
  when it lies in $B_r$, and a permanent cycle when it lies in $Z_\infty$.
  A page class survives to $E_\infty$ when it has a permanent representative;
  it may still represent zero if that representative lies in $B_\infty$.
\end{definition}

% SOURCE: aimpaper/main.tex:253-263, standing convergence assumptions.
\begin{definition}[Strong convergence with detection]
  \label{def:strong-convergence-detection}
  \uses{def:core-filtration,def:complete-separated-filtration,def:chosen-page-presentation}
  \notready
  A strongly convergent elementwise spectral sequence for a graded abutment
  $A_*$ consists of a complete, exhaustive, degreewise separated filtration
  $F^sA_n$ and compatible isomorphisms
  $E_\infty^{s,s+n}\cong F^sA_n/F^{s+1}A_n$.  A class
  $x\in E_\infty^{s,s+n}$ detects $\alpha\in A_n$ when $\alpha\in F^sA_n$
  maps to $x$ in this quotient.
\end{definition}

% SOURCE: KIP126/Core/SpectralSequence/FilteredComplex.lean design note; Mathlib/Algebra/Homology/HomotopyCategory/SpectralObject.lean.
\begin{definition}[Filtered-complex to spectral-object adapter]
  \label{def:filtered-complex-spectral-object-adapter}
  \uses{def:core-filtered-chain-complex,def:complete-separated-filtration,def:mathlib-f2-module-category,thm:mathlib-module-category-abelian,def:mathlib-mapping-cone-spectral-object}
  \notready
  Reindex a filtered chain complex as a cochain complex and regard its
  decreasing filtration inclusions as a functor from the opposite filtration
  order to $\mathbb F_2$-module-valued cochain complexes.  Extend the integer
  filtration to the required endpoint indices using its exhaustive and
  complete data, then precompose Mathlib's mapping-cone spectral object with
  this functor.  The adapter records the chain/cochain sign, filtration-order,
  and endpoint conventions explicitly.
\end{definition}

% SOURCE: aimpaper/main.tex:253-271, Definition def:ess; standard filtered-complex construction used implicitly.
\begin{theorem}[Filtered complex produces a spectral sequence]
  \label{thm:filtered-complex-to-spectral-sequence}
  \uses{def:filtered-complex-spectral-object-adapter,def:mathlib-spectral-sequence-data-core,def:mathlib-has-spectral-sequence,def:mathlib-spectral-object-spectral-sequence,def:mathlib-spectral-object-cycles,def:mathlib-spectral-object-opcycles,def:chosen-page-presentation,def:complete-separated-filtration,def:ss-page-level}
  \notready
  A decreasingly filtered chain complex of $\mathbb F_2$-vector spaces whose
  differential preserves filtration has a functorial spectral sequence with
  $E_0$ the associated-graded complex.  Under exhaustive, complete,
  separated, and degreewise regular convergence data, its $E_\infty$ page is
  the associated graded of filtered homology.  The construction supplies the
  chosen $Z_r/B_r$ presentation and is natural for filtration-preserving chain
  maps.
\end{theorem}

\begin{proof}
  Apply the adapter, define the page recipe for the chosen $E_0$ convention,
  and discharge Mathlib's endpoint-vanishing
  \texttt{HasSpectralSequence} fields from regularity.  Invoke Mathlib's
  spectral-object construction rather than defining a second spectral
  sequence.  Identify its categorical cycles and opcycles with the filtration
  kernels and cokernels and hence with $Z_r$ and $B_r$.  Finally use
  completeness, separatedness, and degreewise regularity to eliminate the
  Boardman obstruction and identify the limiting filtration with filtered
  homology.
\end{proof}

% SOURCE: aimpaper/main.tex:461-631 and 1218-1235, maps with shifted degrees.
\begin{definition}[Reindexed spectral-sequence morphism]
  \label{def:reindexed-ss-morphism}
  \uses{def:mathlib-graded-comap,def:mathlib-complex-shape-up,def:mathlib-spectral-sequence-hom,def:adams-bidegree,def:synthetic-tridegree}
  \notready
  A reindexed morphism specifies a map of page index types, a map of grading
  indices, page maps, and proofs that the page maps commute with differentials
  after the stated degree translation and with page passage.  The interface
  must cover shifts $(s,t)\mapsto(s+k,t+k+d)$ and the forgetful map
  $(s,t,w)\mapsto(s,t)$ used in classical--synthetic comparison.
\end{definition}

% SOURCE: aimpaper/main.tex:138-150, 159-162, 1606-1753.
\begin{definition}[Multiplicative spectral sequence]
  \label{def:multiplicative-ss}
  \uses{def:chosen-page-presentation}
  \notready
  A multiplicative spectral sequence has unital associative products on its
  pages, compatible page maps, additive bidegrees, and a Leibniz rule for page
  differentials with the chosen sign convention.  Its convergence data also
  identifies the $E_\infty$ product with the associated-graded product on the
  abutment.  This interface is required before expressions such as $h_j^2$ or
  $h_0x$ are meaningful.
\end{definition}

% SOURCE: aimpaper/main.tex:253-263 and all later elementwise arguments.
\begin{proposition}[Categorical and elementwise presentations agree]
  \label{prop:categorical-elementwise-comparison}
  \uses{def:mathlib-spectral-sequence,def:mathlib-short-complex-exact,thm:mathlib-exact-iff-image-kernel,def:mathlib-spectral-object-cycles,def:mathlib-spectral-object-opcycles,def:chosen-page-presentation,thm:filtered-complex-to-spectral-sequence}
  \notready
  For a spectral sequence in graded $\mathbb F_2$-vector spaces constructed
  from a chosen exact couple or filtered complex, the Mathlib page objects are
  naturally isomorphic to the quotients $Z_r/B_r$, and the Mathlib page
  differential corresponds to the elementwise differential.
\end{proposition}

\begin{proof}
  Construct the isomorphism on the initial page, identify each next Mathlib
  page with homology, and induct on $r$.  At the inductive step, the kernel and
  image description of homology gives precisely $Z_{r+1}/B_{r+1}$; naturality
  supplies compatibility with page passage.  The construction must be carried
  out for the concrete classical and synthetic presentations, not postulated
  for an arbitrary spectral sequence.
\end{proof}
